\subsection{Bart\'ok et al.\ (2018): general-purpose GAP for silicon}
\label{sec:appendix-example-bartok2018si}

\paragraph{Q1 -- Bibliographic identification.}
Bart\'ok, Kermode, Bernstein, and Cs\'anyi train a single
general-purpose Gaussian Approximation Potential covering crystalline,
liquid, amorphous, surface, defect, and crack-tip silicon
configurations. Published in \emph{Phys.\ Rev.\ X} 8, 041048 (2018);
arXiv:1805.01568.
\lstinputlisting[language=turtle]{artifacts/kg/listings/bartok2018si/q01-bibliographic.ttl}

\paragraph{Q2 -- Material system.}
Bulk silicon is encoded as a single $\MaterialSystemC$; the broad
configurational coverage (diamond, $\beta$-Sn, simple-hexagonal,
hex-diamond, bcc, bc8, fcc, hcp, st12, liquid, amorphous, surfaces,
point defects, crack tips) is recorded as
$\dlRole{materialClass}$ text and as the eight sampling strategies
in Q6.
\lstinputlisting[language=turtle]{artifacts/kg/listings/bartok2018si/q02-material.ttl}

\paragraph{Q3 -- Reference calculation method.}
Reference data come from DFT calculations.
\lstinputlisting[language=turtle]{artifacts/kg/listings/bartok2018si/q03-reference-method-type.ttl}

\paragraph{Q4 -- Reference settings.}
DFT is performed with CASTEP using the PW91 functional, ultrasoft
pseudopotentials, a 250\,eV plane-wave cutoff, and Monkhorst--Pack
k-point meshes with 0.03\,\AA$^{-1}$ spacing for production (with
0.05\,eV electronic smearing); 0.015\,\AA$^{-1}$ for $E(V)$ testing
curves; 0.07\,\AA$^{-1}$ for amorphous re-optimisation.
\lstinputlisting[language=turtle]{artifacts/kg/listings/bartok2018si/q04-reference-settings.ttl}

\paragraph{Q5 -- Training dataset.}
The fitting database has 2{,}475 unit cells contributing 171{,}815
atomic environments and 531{,}710 pieces of electronic-structure data
(energies, forces, virials). Provenance is \emph{in-house} (the
database was assembled by the authors specifically for this work).
\lstinputlisting[language=turtle]{artifacts/kg/listings/bartok2018si/q05-dataset.ttl}

\paragraph{Q6 -- Sampling strategies.}
Eight strategies were combined: crystal-phase enumeration with strain
perturbations; surface enumeration including reconstructions and
decohesion paths; vacancy/divacancy/interstitial point-defect
enumeration; liquid-phase MD; quenched-melt amorphous trajectories;
crack-tip geometries with bond breaking; screw-dislocation cores;
and a single isolated-atom anchor for the cohesive-energy zero.
\lstinputlisting[language=turtle]{artifacts/kg/listings/bartok2018si/q06-sampling.ttl}

\paragraph{Q7 -- MLIP method, components, implementation.}
The method combines a cubic-spline pair-potential repulsive term with
a many-body SOAP-kernel Gaussian-process regression. Training uses
sparse GPR with CUR-selected representative environments. The
implementation is QUIP with the GAP plug-in.
\lstinputlisting[language=turtle]{artifacts/kg/listings/bartok2018si/q07-method.ttl}

\paragraph{Q8 -- Hyperparameter settings.}
Reported settings: SOAP cutoff $r_c=5.0$\,\AA, atomic-Gaussian width
$\sigma_{\text{atom}}=0.5$\,\AA, kernel polynomial power $\zeta=4$,
energy-scale parameter $\delta=3$\,eV, and $M=9{,}000$ representative
environments. Per-class regularisations $\sigma_E$, $\sigma_F$,
$\sigma_V$ are also tuned but not enumerated as
$\HyperparameterSettingC$ instances.
\lstinputlisting[language=turtle]{artifacts/kg/listings/bartok2018si/q08-hyperparameter-settings.ttl}

\paragraph{Q9 -- MLIP run and trained model.}
A single GPR fit produces the GAP-6 silicon potential.
\lstinputlisting[language=turtle]{artifacts/kg/listings/bartok2018si/q09-run-and-model.ttl}

\paragraph{Q10 -- Benchmark results.}
The headline accuracy figure is the median absolute force-component
error on the held-out testing set (grain-boundary, di-interstitial,
GSF-path, amorphous configurations): 0.025\,eV/\AA. The paper
reports many additional per-property and per-phase comparisons (lattice
constants, elastic moduli, surface energies, defect formation
energies, EOS curves) but in derived form rather than as separate
$\BenchmarkResultC$ instances.
\lstinputlisting[language=turtle]{artifacts/kg/listings/bartok2018si/q10-benchmark-results.ttl}

\paragraph{Q11 -- Computational resources.}
\emph{Not reported.}
\lstinputlisting[language=turtle]{artifacts/kg/listings/bartok2018si/q11-resources.ttl}

\paragraph{Q12 -- Gaps.}
Beyond Q11, the paper does not record explicit values for: the
per-class force/virial regularisations $\sigma_F$, $\sigma_V$ (only
their existence and qualitative role); the SOAP $l_{\max}$ and
$n_{\max}$ basis indices that are referenced under
$\dlRole{hasHyperparameter}$ but lack
$\HyperparameterSettingC$ instances; and any $\dlRole{frozenCore}$
treatment beyond the ultrasoft pseudopotential. The paper also
sketches several downstream applications (crack tip propagation,
amorphous-Si melt-quench MD, recrystallisation kinetics) which are
not encoded as separate $\dlConcept{MLIPRun}$ instances.
